% !TEX program = xelatex
% 2027 新领军学科能力二试 · 综合 · 生物（预测训练卷）
\documentclass[cjk=noto, seal=first]{mojing-exam}

\examtitle{2027 新领军学科能力测试 · 预测训练卷}
\secret{训练用卷 ★ 非官方试题}

\begin{document}

\begin{paper}[time={建议 1.5 小时（训练设置）}, full=100]{二试 · 综合 · 生物}

\begin{instructions}
  \item 官方只公布二试综合六选二的科目名称，没有公布生物的考纲、题型与分值。本卷按“高中生物核心内容 + 实验推理与定量分析”编写，
        全部安排为训练设置。
  \item 本卷共 4 题，每题 25 分，满分 100 分。
\end{instructions}

\begin{gradingnotes}
  \item 本评分标准为训练用建议标准。概率可用分数或小数表示；分值为累计得分。
\end{gradingnotes}

\scoretable{1,2,3,4}

\begin{questions}

% =============================================================================
\question[25]
（群体遗传）某常染色体隐性遗传病在一个大群体中的发病率为 $\frac1{10\,000}$。设该群体处于遗传平衡（哈迪—温伯格平衡）。
\begin{enumerate}
  \item 求致病基因频率 $q$，以及表现正常的人中携带者所占的比例；
  \item 两个表现正常、彼此无亲缘关系的人婚配，求孩子患病的概率；
  \item 某表现正常的人，其同胞（兄弟姐妹）中有一人患病，父母均正常。此人与群体中一个表现正常、无亲缘关系的人婚配，求孩子患病的概率；
  \item 说明近亲婚配为什么会提高隐性遗传病的发病风险。
\end{enumerate}
\answerspace{2cm}
\begin{solution}
(1)~$q^2=10^{-4}$，$q=0.01$，$p=0.99$。携带者频率 $2pq=0.0198$，在表现正常者中占 $\frac{2pq}{p^2+2pq}=\frac{2q}{1+q}\approx0.0198\approx\frac1{50}$。\score{7}

(2)~$\bigl(\frac{2q}{1+q}\bigr)^2\times\frac14\approx\frac{1}{50.5^2\times4}\approx9.8\times10^{-5}$，约为 $\frac1{10\,000}$。\score{12}

(3)~父母均为携带者，此人表现正常，是携带者的概率为 $\frac23$；配偶为携带者的概率约 $\frac1{50}$。孩子患病概率
$\frac23\times\frac{1}{50}\times\frac14\approx3.3\times10^{-3}$，约为 (2) 的 34 倍。\score{20}

(4)~近亲双方由共同祖先继承同一个致病基因的概率远大于群体中两个随机个体同为携带者的概率，使隐性纯合子出现的机会显著增加。\score{25}
\end{solution}

% =============================================================================
\question[25]
（连锁与交换）用某植物的纯合品系杂交得到 $F_1$（基因型 AaBb），再用 $F_1$ 与 aabb 测交，子代表型及数量如下：
\begin{center}
\begin{tabular}{cccc}
  \toprule
  AaBb & aabb & Aabb & aaBb \\
  \midrule
  420 & 410 & 90 & 80 \\
  \bottomrule
\end{tabular}
\end{center}
\begin{enumerate}
  \item 判断两对基因是否独立遗传，说明理由；写出 $F_1$ 中两对基因在染色体上的排列方式；
  \item 求 A、B 之间的重组频率与图距；
  \item 另一对基因 C/c 与 A、B 位于同一条染色体上，测得 A—C 的重组频率为 7\%，B—C 的重组频率为 10\%。确定三对基因的排列顺序；
        若无交叉干扰，A 与 B 之间同时发生两次交换（双交换）的子代约占多少？
\end{enumerate}
\answerspace{2cm}
\begin{solution}
(1)~若独立遗传，四种表型应约为 $1:1:1:1$（各约 250），实际亲本型（AaBb、aabb）远多于重组型，说明两对基因连锁。
亲本型为 AB 与 ab，$F_1$ 中 A 与 B 位于同一条染色体上（AB/ab，相引相）。\score{8}

(2)~重组频率 $\frac{90+80}{1000}=17\%$，图距约为 17 cM。\score{14}

(3)~$7+10=17$，C 位于 A、B 之间，顺序为 A—C—B（或 B—C—A）。\score{20}
无干扰时双交换的概率约为 $0.07\times0.10=0.7\%$。\score{25}
\end{solution}

% =============================================================================
\question[25]
（DNA 复制方式的实验推理）把大肠杆菌放在只含 $^{15}$N 的培养基中培养多代，使其 DNA 全部含 $^{15}$N（“重”），然后转移到只含 $^{14}$N
的培养基中繁殖。每繁殖一代取样，提取 DNA 进行密度梯度离心，根据条带位置区分“重”（两条链均含 $^{15}$N）、“中”（一条链 $^{15}$N、一条链
$^{14}$N）、“轻”（两条链均含 $^{14}$N）。关于 DNA 复制方式有三种假说：全保留复制、半保留复制、弥散复制（新旧材料在每条链中交错分布）。
\begin{enumerate}
  \item 分别预测三种假说下第 1 代、第 2 代的离心条带；
  \item 实验结果：第 1 代只有一条“中”带；第 2 代有“中”“轻”两条带，且量相等。哪个假说成立？为什么只看第 1 代还不能下结论？
  \item 按成立的假说，第 $n$ 代（$n\ge1$）中“中”DNA 所占的比例是多少？
\end{enumerate}
\answerspace{2cm}
\begin{solution}
(1)~全保留：第 1 代“重”“轻”各一条带，第 2 代“重”$:$“轻”$=1:3$。半保留：第 1 代只有“中”带；第 2 代“中”“轻”各半。
弥散：第 1 代只有一条中间密度带；第 2 代仍只有一条带，位置在“中”与“轻”之间。\score{12}

(2)~半保留复制。第 1 代结果已排除全保留，但半保留与弥散在第 1 代都只出现一条“中”密度带，无法区分；第 2 代出现两条分开的带，排除弥散。\score{20}

(3)~原有的两条 $^{15}$N 链始终各在一个 DNA 分子中，第 $n$ 代共 $2^n$ 个分子，“中”DNA 占 $\dfrac{2}{2^n}=2^{1-n}$。\score{25}
\end{solution}

% =============================================================================
\question[25]
（种群增长与捕捞）某鱼类种群数量 $N(t)$ 满足逻辑斯谛模型 $\dfrac{\dd N}{\dd t}=rN\Bigl(1-\dfrac NK\Bigr)$，$r>0$ 为内禀增长率，$K$ 为环境容纳量。
现以恒定速率 $H$ 捕捞，模型变为 $\dfrac{\dd N}{\dd t}=rN\Bigl(1-\dfrac NK\Bigr)-H$。
\begin{enumerate}
  \item 不捕捞时，种群数量为多少时增长最快？
  \item 求使种群能长期维持的最大捕捞速率 $H_{\max}$（最大持续产量）；当 $0<H<H_{\max}$ 时，求平衡点并判断哪个是稳定的；
  \item 取 $r=0.4$ /年，$K=1.0\times10^4$ 尾，求 $H_{\max}$。若实际捕捞速率比 $H_{\max}$ 稍大，种群将怎样变化？这对渔业管理有什么启示？
\end{enumerate}
\answerspace{2cm}
\begin{solution}
(1)~$g(N)=rN(1-N/K)$ 在 $N=K/2$ 时取最大值 $rK/4$，此时增长最快。\score{5}

(2)~平衡点满足 $g(N)=H$。由 (1)，$H\le rK/4$ 时才有平衡点，故 $H_{\max}=\frac{rK}4$。\score{11}
$0<H<H_{\max}$ 时有两个平衡点 $N_\pm=\frac K2\Bigl(1\pm\sqrt{1-\frac{4H}{rK}}\Bigr)$。在 $N_+$ 附近，$N$ 略大则 $g(N)<H$ 使 $N$ 减小，$N$ 略小则增大，
故 $N_+$ 稳定；$N_-$ 不稳定，种群若低于 $N_-$ 将走向灭绝。\score{19}

(3)~$H_{\max}=\frac{0.4\times10^4}{4}=1000$ 尾/年。若 $H>H_{\max}$，恒有 $\frac{\dd N}{\dd t}<0$，种群将在有限时间内崩溃。
启示：按理论最大持续产量捕捞没有安全余量，估计误差或环境波动都可能导致崩溃；应低于 $H_{\max}$ 捕捞，或按种群数量比例（而非固定数量）捕捞。\score{25}
\end{solution}

\end{questions}
\end{paper}
\end{document}
